\chapter{Phân tích thiết kế hệ thống}

Chương này đi sâu vào phân tích các yêu cầu chức năng và phi chức năng của hệ thống Synergit, được mô hình hoá chi tiết qua lược đồ Usecase và danh sách tác nhân. Đồng thời, chương trình bày thiết kế \textit{kiến trúc Microservices} với 4 dịch vụ độc lập, thiết kế cơ sở dữ liệu quan hệ (ERD) cùng đặc tả cấu trúc các bảng dữ liệu, và các nghiệp vụ đặc thù như giao thức Git Smart HTTP, máy trạng thái pull request, máy trạng thái issue.

\section{Xác định yêu cầu}

\subsection{Khảo sát hiện trạng các nền tảng quản lý mã nguồn}
Trước khi xây dựng hệ thống, nhóm đã khảo sát các nền tảng quản lý mã nguồn phổ biến hiện nay. Bảng~\ref{tab:khaosat} so sánh các nền tảng tiêu biểu trên những tiêu chí quan trọng đối với một đội phát triển nội bộ tại Việt Nam.

\renewcommand{\arraystretch}{1.3}
\begin{longtable}{|>{\raggedright\arraybackslash}m{2.9cm}|>{\centering\arraybackslash}m{2.6cm}|>{\centering\arraybackslash}m{2.4cm}|>{\centering\arraybackslash}m{2.4cm}|>{\centering\arraybackslash}m{2.2cm}|}
\caption{Bảng khảo sát các nền tảng quản lý mã nguồn}\label{tab:khaosat}\\
\hline\rowcolor{sgCream}
\textbf{Tiêu chí} & \textbf{GitHub.com (SaaS)} & \textbf{GitLab Self-hosted} & \textbf{Gitea Self-hosted} & \textbf{Synergit}\\
\hline\endfirsthead
\hline\rowcolor{sgCream}\textbf{Tiêu chí} & \textbf{GitHub.com (SaaS)} & \textbf{GitLab Self-hosted} & \textbf{Gitea Self-hosted} & \textbf{Synergit}\\\hline\endhead
Tự triển khai (self-hosted) & Không & Có & Có & \textbf{Có}\\\hline
Sở hữu dữ liệu \& mã nguồn & Không & Có & Có & \textbf{Có}\\\hline
Quy trình PR \& Issue đầy đủ & Có & Có & Có & \textbf{Có}\\\hline
Insights / Analytics tích hợp & Có (nâng cao) & Có (nâng cao) & Hạn chế & \textbf{Có (cơ bản)}\\\hline
Kiến trúc microservices rõ ràng & N/A (kín) & Một phần & Không (monolith) & \textbf{Có (4 dịch vụ)}\\\hline
Mã nguồn mở (open source) & Không & Có (CE) & Có & N/A (đồ án)\\\hline
Yêu cầu tài nguyên triển khai & N/A & Cao (Ruby + PG + Redis + ElasticSearch+\ldots) & Thấp & \textbf{Trung bình}\\\hline
Mức độ phức tạp vận hành & Không (do nhà cung cấp) & Cao & Thấp & \textbf{Trung bình}\\\hline
Mục đích chính & Sản phẩm SaaS & Sản phẩm DevOps đầy đủ & Git server gọn nhẹ & \textbf{Đồ án mô phỏng kiến trúc}\\\hline
\end{longtable}

\noindent\textbf{Nhận xét:} GitHub.com là tiêu chuẩn vàng nhưng là dịch vụ kín và không tự triển khai được; GitLab CE rất mạnh nhưng quá ``nặng'' cho nhu cầu nội bộ vừa và nhỏ; Gitea gọn nhẹ nhưng chỉ là một monolith Go đơn giản, không có insights và không phải mô hình tham khảo cho microservices. Synergit hướng tới dung hoà: cung cấp các tính năng cộng tác cốt lõi giống GitHub (PR, issue, insights) trong một kiến trúc microservices rõ ràng làm mô hình tham khảo, đồng thời vẫn nhẹ đủ để chạy được trên một máy phát triển.

\subsection{Những vấn đề còn tồn tại}
\begin{itemize}
\item \textbf{Phụ thuộc dịch vụ nước ngoài:} Mã nguồn nội bộ (đặc biệt là dự án nhạy cảm) đặt trên GitHub/Bitbucket gây lo ngại về chủ quyền dữ liệu.
\item \textbf{Chi phí sử dụng theo người dùng:} Các nền tảng SaaS tính phí theo seat; chi phí tăng tuyến tính khi đội phát triển mở rộng.
\item \textbf{Nền tảng tự triển khai sẵn quá phức tạp hoặc quá đơn giản:} GitLab self-hosted yêu cầu tài nguyên cao và kiến trúc phức tạp; Gitea đơn giản nhưng thiếu các tính năng phân tích và là một monolith chặt chẽ -- không phải mô hình tham khảo cho microservices.
\item \textbf{Thiếu mô hình tham khảo kiến trúc:} Sinh viên/kỹ sư trẻ ít có cơ hội tiếp cận mã nguồn của một hệ thống quản lý mã nguồn được tách dịch vụ rõ ràng theo DDD.
\end{itemize}

\subsection{Vấn đề đề tài tập trung giải quyết}
\begin{itemize}
\item \textbf{Tự triển khai \& sở hữu dữ liệu:} Synergit là phần mềm có thể tự triển khai trên máy chủ nội bộ; toàn bộ mã nguồn và dữ liệu nằm trong tầm kiểm soát của tổ chức.
\item \textbf{Đầy đủ tính năng cộng tác cốt lõi:} Pull request, code review qua diff, issue tracker, label, assignee, comment, event timeline, insights, contribution graph.
\item \textbf{Kiến trúc Microservices rõ ràng theo DDD:} Mô hình tham khảo cho việc tách dịch vụ theo bounded context; tận dụng Clean Architecture để tách dịch vụ trở thành thao tác kỹ thuật minh bạch (thay adapter).
\item \textbf{Đơn giản trong vận hành:} Chỉ 4 binary Go + 1 Postgres + frontend tĩnh; có thể chạy bằng Docker Compose trên một máy phát triển.
\end{itemize}

\section{Phân tích yêu cầu}

\subsection{Yêu cầu chức năng}
\textbf{Nhóm chức năng cho người dùng đăng ký (Authenticated User):}
\begin{itemize}
\item Quản lý tài khoản: đăng ký, đăng nhập, đổi tên đăng nhập, đổi mật khẩu, đăng xuất.
\item Quản lý hồ sơ cá nhân: ảnh đại diện, mô tả; xem trang profile của bản thân.
\item Tạo và quản lý repository: tạo mới (với tuỳ chọn khởi tạo README/.gitignore/LICENSE), đổi tên, đổi visibility, xoá.
\item Push/clone/pull mã nguồn qua Git Smart HTTP cho repository sở hữu hoặc được cấp quyền.
\item Duyệt repository: cây thư mục, nội dung file, lịch sử commit, chi tiết commit (diff).
\item Quản lý branch: tạo từ branch khác, đổi tên, xoá.
\item Commit thay đổi qua giao diện web: chỉnh sửa file đơn lẻ hoặc nhiều file cùng lúc.
\item Tạo và quản lý pull request: so sánh branch, tạo PR, merge, revert, đóng/mở lại, gán label, gán assignee.
\item Giải quyết xung đột (merge conflict) qua giao diện web.
\item Tạo và quản lý issue: tạo, gán label/assignee, comment, đóng (kèm lý do)/mở lại.
\item Quản lý label của repository (tạo, sửa, xoá, áp dụng cho issue và PR).
\item Quản lý collaborator: thêm/gỡ thành viên với vai trò OWNER/MAINTAINER/WRITE.
\item Đánh dấu sao (star) repository; xem danh sách repository đã star.
\item Xem Repo Insights (commit trend, top contributor, language breakdown\ldots) và Profile Activity (contribution graph 365 ngày).
\end{itemize}

\textbf{Nhóm chức năng cho khách (Guest -- chưa đăng nhập):}
\begin{itemize}
\item Đăng ký tài khoản mới, đăng nhập tài khoản đã có.
\item Xem các repository công khai (PUBLIC) của người dùng khác qua URL.
\item Clone hoặc pull repository công khai qua Git Smart HTTP (không cần xác thực).
\item Xem trang profile công khai của người dùng khác.
\end{itemize}

\subsection{Yêu cầu phi chức năng}
Hệ thống cần đáp ứng các tiêu chuẩn kỹ thuật sau:

\subsubsection{Hiệu năng (Performance)}
\begin{itemize}
\item Các trang nội dung chính (cây thư mục, danh sách commit, danh sách PR/issue) tải xong trong vòng dưới 2 giây ở quy mô nhỏ (< 100 commit/PR/issue).
\item Các thao tác Git \texttt{clone}/\texttt{push}/\texttt{pull} thực hiện trực tiếp qua go-git, không phải shell-out, tận dụng tối đa concurrency của Go.
\item Tách các tác vụ phân tích nặng (recompute insights) ra dịch vụ riêng để không ảnh hưởng tới luồng nghiệp vụ chính.
\end{itemize}

\subsubsection{Bảo mật (Security)}
\begin{itemize}
\item Xác thực bằng JWT (HMAC-SHA256, hạn 7 ngày), token gắn vào header \texttt{Authorization: Bearer ...}.
\item Mật khẩu băm một chiều bằng bcrypt; không lưu mật khẩu thô.
\item Phân quyền theo vai trò collaborator (OWNER/MAINTAINER/WRITE); kiểm tra ở tầng usecase.
\item Visibility của repository: PUBLIC ai cũng xem được (kể cả guest); PRIVATE chỉ owner và collaborator xem được.
\item Tất cả dịch vụ backend cùng dùng một secret JWT, tránh việc gọi xác thực trung tâm tăng độ trễ.
\item Validate dữ liệu đầu vào ở tầng handler; toàn bộ giao tiếp qua HTTPS khi triển khai.
\end{itemize}

\subsubsection{Khả năng mở rộng (Scalability)}
\begin{itemize}
\item Kiến trúc microservices với 4 dịch vụ độc lập, mỗi dịch vụ có thể nhân bản (scale horizontal) độc lập.
\item Schema-per-service đảm bảo độc lập dữ liệu, tránh ràng buộc cứng giữa các dịch vụ.
\item Clean Architecture cho phép thay đổi adapter (PostgreSQL $\to$ Postgres khác, hệ thống tệp $\to$ object storage) mà không ảnh hưởng usecase.
\item Đóng gói bằng Docker để dễ nhân bản (scale) khi tải tăng.
\item Tương thích các trình duyệt phổ biến (Chrome, Firefox, Edge, Safari).
\end{itemize}

\subsubsection{Trải nghiệm người dùng (Usability \& UX)}
\begin{itemize}
\item Giao diện lấy cảm hứng trực tiếp từ GitHub: layout, hệ màu, thuật ngữ, icon (Octicons) -- người dùng quen với GitHub có thể sử dụng Synergit ngay lập tức mà không cần học lại.
\item Code browser hỗ trợ tô màu cú pháp; diff viewer hiển thị thêm/bớt theo dòng kiểu unified diff.
\item Contribution graph 365 ngày được vẽ chính xác kiểu GitHub (heatmap 7 hàng x ~53 cột).
\item Thông báo lỗi rõ ràng, có hướng dẫn nhập liệu (validation) ngay trên biểu mẫu.
\end{itemize}

\section{Mô hình hoá chức năng hệ thống (Usecase)}

\subsection{Danh sách tác nhân (Actors)}
\renewcommand{\arraystretch}{1.25}
\begin{longtable}{|c|p{3.6cm}|p{9.3cm}|}
\caption{Danh sách tác nhân của hệ thống}\label{tab:actors}\\
\hline\rowcolor{sgCream}\textbf{STT} & \textbf{Tác nhân} & \textbf{Mô tả}\\\hline\endfirsthead
\hline\rowcolor{sgCream}\textbf{STT} & \textbf{Tác nhân} & \textbf{Mô tả}\\\hline\endhead
1 & Khách (Guest) & Người dùng chưa đăng nhập; xem repository công khai, profile công khai, clone qua Git.\\\hline
2 & Người dùng (Authenticated User) & Người dùng đã đăng ký; tạo repository riêng, push/pull, tạo PR/issue.\\\hline
3 & Chủ sở hữu (Owner) & Người dùng tạo repository; toàn quyền: quản lý collaborator, đổi visibility, xoá repo, merge PR.\\\hline
4 & Maintainer & Cộng tác viên với quyền cao: thêm/gỡ collaborator (trừ OWNER), merge PR, quản lý label.\\\hline
5 & Writer & Cộng tác viên với quyền ghi: push, commit, tạo PR/issue, comment.\\\hline
\end{longtable}

\subsection{Sơ đồ Usecase tổng quát}
Sơ đồ Usecase tổng quát mô tả tương tác giữa các tác nhân và các nhóm chức năng chính của hệ thống được trình bày trong Hình~\ref{fig:usecase}.

% ====================== Sơ đồ Usecase tổng quát Synergit (TikZ) ============
\begin{figure}[H]\centering
\resizebox{\linewidth}{!}{%
\begin{tikzpicture}[
  font=\sffamily,
  actor/.pic={
    \draw[thick] (0,0.32)  circle [radius=0.14];
    \draw[thick] (0,0.18) -- (0,-0.20);
    \draw[thick] (-0.22,0.02) -- (0.22,0.02);
    \draw[thick] (0,-0.20) -- (-0.18,-0.50);
    \draw[thick] (0,-0.20) -- (0.18,-0.50);
  },
  albl/.style={font=\scriptsize\bfseries, align=center},
  uc/.style={draw=sgBlue, fill=sgCream, ellipse, align=center,
             font=\scriptsize, inner sep=1pt, minimum width=3.4cm, minimum height=1.0cm,
             text width=3.0cm},
  assoc/.style={draw=sgLine, line width=0.4pt},
]

% ---------------- Tác nhân (trái) ----------------
\node (guest) at (0,8.0) {}; \pic at (guest) {actor}; \node[albl,below=0.45cm of guest] {Khách\\(Guest)};
\node (user)  at (0,4.0) {}; \pic at (user)  {actor}; \node[albl,below=0.45cm of user]  {Người dùng\\đã đăng ký};

% ---------------- Tác nhân (phải) ----------------
\node (writer) at (14,8.6) {}; \pic at (writer) {actor}; \node[albl,below=0.45cm of writer] {Writer\\(WRITE)};
\node (maint)  at (14,5.6) {}; \pic at (maint)  {actor}; \node[albl,below=0.45cm of maint]  {Maintainer};
\node (owner)  at (14,2.6) {}; \pic at (owner)  {actor}; \node[albl,below=0.45cm of owner]  {Chủ sở hữu\\(Owner)};

% ---------------- Use cases: cột trái (người dùng / khách) ----------------
\node[uc] (u1) at (4.6,8.6)  {Đăng ký / Đăng nhập};
\node[uc] (u2) at (4.6,7.2)  {Xem repo công khai};
\node[uc] (u3) at (4.6,5.8)  {Clone / Pull qua Git};
\node[uc] (u4) at (4.6,4.4)  {Tạo repository};
\node[uc] (u5) at (4.6,3.0)  {Đổi tên đăng nhập};
\node[uc] (u6) at (4.6,1.6)  {Xem Profile Activity};
\node[uc] (u7) at (4.6,0.2)  {Quản lý hồ sơ};

% ---------------- Use cases: cột phải (cộng tác trên repo) ----------------
\node[uc] (r1) at (9.4,9.0)  {Push / Commit qua web};
\node[uc] (r2) at (9.4,7.7)  {Quản lý branch};
\node[uc] (r3) at (9.4,6.4)  {Tạo \& bình luận PR};
\node[uc] (r4) at (9.4,5.1)  {Tạo \& quản lý Issue};
\node[uc] (r5) at (9.4,3.8)  {Merge / Revert PR};
\node[uc] (r6) at (9.4,2.5)  {Resolve conflict};
\node[uc] (r7) at (9.4,1.2)  {Quản lý Collaborator};
\node[uc] (r8) at (9.4,-0.1) {Đổi visibility / Xoá repo};

% ---------------- Ranh giới hệ thống ----------------
\begin{scope}[on background layer]
\node[draw=sgBlue, very thick, rounded corners, fill=sgBlue!3,
      fit=(u1)(u7)(r1)(r8), inner sep=16pt,
      label={[font=\sffamily\bfseries, text=sgBlue]above:HỆ THỐNG SYNERGIT}] {};
\end{scope}

% ---------------- Liên kết tác nhân trái ----------------
\draw[assoc] (guest) -- (u1.west);
\draw[assoc] (guest) -- (u2.west);
\draw[assoc] (guest) -- (u3.west);
\draw[assoc] (user) -- (u4.west);
\draw[assoc] (user) -- (u5.west);
\draw[assoc] (user) -- (u6.west);
\draw[assoc] (user) -- (u7.west);
\draw[assoc] (user) -- (u3.west);

% ---------------- Liên kết tác nhân phải ----------------
% Writer (push, branch, tạo PR/issue/comment)
\draw[assoc] (writer) -- (r1.east);
\draw[assoc] (writer) -- (r2.east);
\draw[assoc] (writer) -- (r3.east);
\draw[assoc] (writer) -- (r4.east);
% Maintainer
\draw[assoc] (maint) -- (r5.east);
\draw[assoc] (maint) -- (r6.east);
\draw[assoc] (maint) -- (r7.east);
% Owner: bao trùm tất cả (chỉ vẽ một số quan trọng để gọn)
\draw[assoc] (owner) -- (r5.east);
\draw[assoc] (owner) -- (r7.east);
\draw[assoc] (owner) -- (r8.east);
\end{tikzpicture}}
\caption{Sơ đồ Usecase tổng quát của hệ thống Synergit}\label{fig:usecase}
\end{figure}


\subsection{Danh sách Usecase}
\renewcommand{\arraystretch}{1.2}
\begin{longtable}{|c|c|p{5.2cm}|p{6.0cm}|}
\caption{Danh sách Usecase chính}\label{tab:uclist}\\
\hline\rowcolor{sgCream}\textbf{STT} & \textbf{ID} & \textbf{Tên usecase} & \textbf{Mô tả}\\\hline\endfirsthead
\hline\rowcolor{sgCream}\textbf{STT} & \textbf{ID} & \textbf{Tên usecase} & \textbf{Mô tả}\\\hline\endhead
1  & UC-01 & Đăng ký tài khoản & Khách tạo tài khoản mới.\\\hline
2  & UC-02 & Đăng nhập & Người dùng đăng nhập bằng tên đăng nhập/mật khẩu.\\\hline
3  & UC-03 & Đổi tên đăng nhập & Người dùng đổi username; hệ thống đổi tên thư mục lưu repo và phát hành JWT mới.\\\hline
4  & UC-04 & Tạo repository & Tạo repo mới với tuỳ chọn README, .gitignore, LICENSE.\\\hline
5  & UC-05 & Quản lý repository & Đổi tên, đổi visibility, xoá repo.\\\hline
6  & UC-06 & Push/Clone/Pull qua Git & Thao tác Git Smart HTTP với repo.\\\hline
7  & UC-07 & Duyệt cây thư mục \& xem file & Xem tree theo branch, xem nội dung file.\\\hline
8  & UC-08 & Xem lịch sử commit & Liệt kê commit theo branch, lọc theo path.\\\hline
9  & UC-09 & Xem chi tiết commit (diff) & Xem diff của một commit theo từng file.\\\hline
10 & UC-10 & Quản lý branch & Tạo, đổi tên, xoá; xem ahead/behind so với mặc định.\\\hline
11 & UC-11 & Commit qua giao diện web & Chỉnh sửa file đơn lẻ hoặc nhiều file rồi commit.\\\hline
12 & UC-12 & So sánh hai branch (compare) & Hiển thị commit và file thay đổi giữa base và head.\\\hline
13 & UC-13 & Tạo Pull Request & Tạo PR từ kết quả compare; nhập title, description.\\\hline
14 & UC-14 & Xem PR \& timeline & Xem chi tiết PR cùng dòng thời gian sự kiện.\\\hline
15 & UC-15 & Merge Pull Request & Merge nếu sạch; chuyển luồng resolve khi có xung đột.\\\hline
16 & UC-16 & Giải quyết xung đột (merge conflict) & Hiển thị file xung đột với marker, người dùng nhập nội dung đã giải quyết, hệ thống commit.\\\hline
17 & UC-17 & Revert Pull Request đã merge & Tạo branch revert + commit revert tự động.\\\hline
18 & UC-18 & Đóng/Mở lại PR & Chuyển trạng thái PR.\\\hline
19 & UC-19 & Tạo Issue & Tạo issue (title, description) cho repo.\\\hline
20 & UC-20 & Quản lý Issue & Comment, gán/bỏ assignee, đóng (kèm lý do)/mở lại.\\\hline
21 & UC-21 & Quản lý Label & Tạo, sửa, xoá label; gán cho issue/PR.\\\hline
22 & UC-22 & Quản lý Collaborator & Thêm/gỡ collaborator với vai trò OWNER/MAINTAINER/WRITE.\\\hline
23 & UC-23 & Star/Unstar repository & Đánh dấu yêu thích; xem starred của bản thân.\\\hline
24 & UC-24 & Xem Repo Insights & Commit trend, contributors, branch activity, language breakdown.\\\hline
25 & UC-25 & Recompute Insights & Người dùng chủ động tính lại insights cho repo.\\\hline
26 & UC-26 & Xem Profile Activity & Contribution graph 365 ngày, activity overview.\\\hline
\end{longtable}

\subsection{Mô tả các Usecase chính}
Phần này đặc tả chi tiết các usecase tiêu biểu, đại diện đầy đủ cho các luồng nghiệp vụ và các tác nhân của hệ thống.

% ====================== Đặc tả các Usecase chính ============================

\paragraph{UC-01: Đăng ký tài khoản}
\begin{ucspec}{uc01}{Đặc tả Usecase Đăng ký tài khoản}
Use Case ID & UC-01 \\\hline
Tên Usecase & Đăng ký tài khoản \\\hline
Mô tả & Khách tạo tài khoản mới bằng tên đăng nhập và mật khẩu. \\\hline
Tác nhân & Khách (Guest) \\\hline
Độ ưu tiên & Cao \\\hline
Tiền điều kiện & Tên đăng nhập và email chưa tồn tại trong hệ thống. \\\hline
Hậu điều kiện & Tài khoản được tạo, người dùng có thể đăng nhập ngay lập tức. \\\hline
Luồng chính &
1. Khách nhập tên đăng nhập, email, mật khẩu (tối thiểu 6 ký tự).\newline
2. Hệ thống kiểm tra hợp lệ (username/email chưa tồn tại, mật khẩu đủ mạnh, email có ký tự ``@'').\newline
3. Hệ thống băm mật khẩu (bcrypt) và lưu bản ghi mới vào \texttt{core.users}.\newline
4. Trả về thông báo thành công; chuyển sang trang đăng nhập (UC-02). \\\hline
Luồng ngoại lệ & A1: Username/email đã tồn tại $\rightarrow$ báo lỗi và đề nghị đăng nhập.\newline A2: Mật khẩu yếu (< 6 ký tự) $\rightarrow$ hiển thị yêu cầu cụ thể. \\\hline
\end{ucspec}

\paragraph{UC-02: Đăng nhập}
\begin{ucspec}{uc02}{Đặc tả Usecase Đăng nhập}
Use Case ID & UC-02 \\\hline
Tên Usecase & Đăng nhập \\\hline
Mô tả & Người dùng đăng nhập bằng tên đăng nhập/email và mật khẩu để nhận JWT. \\\hline
Tác nhân & Người dùng đã đăng ký \\\hline
Độ ưu tiên & Cao \\\hline
Tiền điều kiện & Tài khoản đã tồn tại. \\\hline
Hậu điều kiện & Người dùng nhận được JWT, lưu trên client; được chuyển tới dashboard. \\\hline
Luồng chính &
1. Người dùng nhập username (hoặc email) và mật khẩu.\newline
2. Hệ thống tìm bản ghi trong \texttt{core.users} và so khớp \texttt{password\_hash} bằng bcrypt.\newline
3. Nếu hợp lệ, sinh JWT (HMAC-SHA256, hạn 7 ngày) chứa \texttt{user\_id} và \texttt{username}.\newline
4. Trả token cho client; client lưu vào \texttt{localStorage} và đính kèm header \texttt{Authorization: Bearer ...} cho mọi request sau đó.\newline
5. Chuyển hướng tới dashboard repository. \\\hline
Luồng ngoại lệ & A1: Sai username/mật khẩu $\rightarrow$ báo lỗi chung (không tiết lộ trường nào sai). \\\hline
\end{ucspec}

\paragraph{UC-03: Đổi tên đăng nhập}
\begin{ucspec}{uc03}{Đặc tả Usecase Đổi tên đăng nhập}
Use Case ID & UC-03 \\\hline
Tên Usecase & Đổi tên đăng nhập \\\hline
Mô tả & Người dùng đổi username; hệ thống đổi tên thư mục lưu repo và phát hành JWT mới. \\\hline
Tác nhân & Người dùng \\\hline
Độ ưu tiên & Trung bình \\\hline
Tiền điều kiện & Người dùng đã đăng nhập; username mới chưa tồn tại. \\\hline
Hậu điều kiện & Username được cập nhật ở DB, hệ thống tệp, và trong JWT. \\\hline
Luồng chính &
1. Người dùng vào trang Settings, nhập username mới.\newline
2. Core kiểm tra username mới chưa tồn tại trong \texttt{core.users}.\newline
3. Cập nhật \texttt{core.users.username} cho user.\newline
4. Cập nhật cột \texttt{path} của các bản ghi trong \texttt{core.repositories} (thay tiền tố \texttt{<old>/} bằng \texttt{<new>/}).\newline
5. Gọi GitStorage qua HTTP (\texttt{POST /git/users/:old/rename}) để đổi tên thư mục \texttt{<gitRoot>\textbackslash<old>\textbackslash}~$\rightarrow$~\texttt{<gitRoot>\textbackslash<new>\textbackslash}.\newline
6. Phát hành JWT mới chứa username mới và trả về client.\newline
7. Client lưu JWT mới vào \texttt{localStorage} và reload trang. \\\hline
Luồng ngoại lệ & A1: Username mới đã tồn tại $\rightarrow$ HTTP 409 Conflict.\newline A2: Username quá ngắn (< 3 ký tự) $\rightarrow$ HTTP 400 Bad Request. \\\hline
\end{ucspec}

\paragraph{UC-04: Tạo repository}
\begin{ucspec}{uc04}{Đặc tả Usecase Tạo repository}
Use Case ID & UC-04 \\\hline
Tên Usecase & Tạo repository \\\hline
Mô tả & Người dùng tạo repository mới với tuỳ chọn khởi tạo README/.gitignore/LICENSE. \\\hline
Tác nhân & Người dùng \\\hline
Độ ưu tiên & Cao \\\hline
Tiền điều kiện & Người dùng đã đăng nhập; repo cùng tên chưa tồn tại trong tài khoản. \\\hline
Hậu điều kiện & Repo được tạo trong DB và trên hệ thống tệp; user là OWNER; nếu chọn khởi tạo, repo có commit đầu tiên. \\\hline
Luồng chính &
1. Người dùng nhập tên repo, mô tả, chọn visibility (PUBLIC/PRIVATE), bật/tắt khởi tạo README và chọn .gitignore/LICENSE template.\newline
2. Core kiểm tra hợp lệ: tên không rỗng, không trùng với repo khác cùng owner.\newline
3. Core gọi GitStorage \texttt{InitBareRepo(<owner>/<name>)} để tạo bare repository trên hệ thống tệp.\newline
4. Nếu chọn khởi tạo, Core gọi GitStorage \texttt{BootstrapRepository} với danh sách file ban đầu (README, .gitignore, LICENSE) và message \texttt{Initial commit}.\newline
5. Core ghi bản ghi vào \texttt{core.repositories} và \texttt{core.repository\_collaborators} (vai trò OWNER).\newline
6. Core seed 9 label mặc định kiểu GitHub vào \texttt{core.labels}.\newline
7. Core trigger Insights tính lần đầu (background, không chặn).\newline
8. Trả về thông tin repo, frontend chuyển tới trang Code của repo. \\\hline
Luồng ngoại lệ & A1: Tên repo trùng $\rightarrow$ HTTP 409.\newline A2: Tên rỗng/không hợp lệ $\rightarrow$ HTTP 400. \\\hline
\end{ucspec}

\paragraph{UC-06: Push/Clone/Pull qua Git Smart HTTP}
\begin{ucspec}{uc06}{Đặc tả Usecase Push/Clone/Pull qua Git Smart HTTP}
Use Case ID & UC-06 \\\hline
Tên Usecase & Push/Clone/Pull qua Git Smart HTTP \\\hline
Mô tả & Người dùng dùng \texttt{git} CLI tương tác với repo qua HTTP. \\\hline
Tác nhân & Người dùng / Khách (cho repo PUBLIC, chỉ đọc) \\\hline
Độ ưu tiên & Cao \\\hline
Tiền điều kiện & Repo tồn tại trên hệ thống tệp; visibility PUBLIC hoặc client có quyền (PRIVATE). \\\hline
Hậu điều kiện & Object được trao đổi đúng giao thức Git Smart HTTP. \\\hline
Luồng chính &
1. Client gọi \texttt{git clone http://gateway/<owner>/<repo>.git} (hoặc \texttt{push}/\texttt{pull}).\newline
2. Gateway nhận, định tuyến tới GitStorage.\newline
3. GitStorage chạy theo nghiệp vụ Smart HTTP:\newline\hspace*{1em}- \texttt{GET .../info/refs?service=git-upload-pack} (hoặc \texttt{git-receive-pack}).\newline\hspace*{1em}- \texttt{POST .../git-upload-pack} hoặc \texttt{git-receive-pack}.\newline
4. GitStorage dùng go-git để parse/sinh packfile và phản hồi.\newline
5. Client áp dụng cập nhật (clone/checkout/push refs). \\\hline
Luồng ngoại lệ & A1: Repo PRIVATE và không có quyền $\rightarrow$ HTTP 401.\newline A2: Repo không tồn tại $\rightarrow$ HTTP 404. \\\hline
\end{ucspec}

\paragraph{UC-07: Duyệt cây thư mục và xem nội dung file}
\begin{ucspec}{uc07}{Đặc tả Usecase Duyệt cây thư mục và xem nội dung file}
Use Case ID & UC-07 \\\hline
Tên Usecase & Duyệt cây thư mục \& xem nội dung file \\\hline
Mô tả & Người dùng duyệt cây thư mục theo branch và xem nội dung file. \\\hline
Tác nhân & Người dùng / Khách (cho repo PUBLIC) \\\hline
Độ ưu tiên & Cao \\\hline
Tiền điều kiện & Repo tồn tại; branch tồn tại. \\\hline
Hậu điều kiện & Hiển thị danh sách entry hoặc nội dung file. \\\hline
Luồng chính &
1. Frontend gọi \texttt{GET /api/v1/repos/:id/tree?path=<path>\&branch=<b>}.\newline
2. Core gọi GitStorage \texttt{GetTree(repoPath, path, branch)} qua HTTP.\newline
3. GitStorage parse tree object qua go-git, trả về danh sách \texttt{RepoFile\{name, path, type\}}.\newline
4. Khi nhấp file, frontend gọi \texttt{GET /api/v1/repos/:id/blob?path=<file>\&branch=<b>}.\newline
5. GitStorage trả về nội dung blob; frontend tô màu cú pháp theo phần mở rộng. \\\hline
Luồng ngoại lệ & A1: Path không tồn tại $\rightarrow$ HTTP 404.\newline A2: Branch không tồn tại $\rightarrow$ HTTP 404. \\\hline
\end{ucspec}

\paragraph{UC-09: Xem chi tiết commit và diff}
\begin{ucspec}{uc09}{Đặc tả Usecase Xem chi tiết commit và diff}
Use Case ID & UC-09 \\\hline
Tên Usecase & Xem chi tiết commit (diff) \\\hline
Mô tả & Người dùng xem nội dung diff của một commit theo từng file. \\\hline
Tác nhân & Người dùng / Khách (cho repo PUBLIC) \\\hline
Độ ưu tiên & Cao \\\hline
Tiền điều kiện & Commit tồn tại trong repo. \\\hline
Hậu điều kiện & Hiển thị danh sách file thay đổi với patch unified diff. \\\hline
Luồng chính &
1. Frontend gọi \texttt{GET /api/v1/repos/:id/commits/:hash} để lấy thông tin commit (author, message, date).\newline
2. Frontend gọi \texttt{GET /api/v1/repos/:id/commits/:hash/diff} để lấy danh sách \texttt{DiffFile}.\newline
3. Core forward sang GitStorage \texttt{GetCommitDiff}.\newline
4. GitStorage so sánh tree của commit với tree của parent (commit gốc thì so với rỗng), sinh patch unified.\newline
5. Frontend hiển thị side-by-side hoặc unified, tô màu xanh (added) và đỏ (deleted) theo dòng. \\\hline
Luồng ngoại lệ & A1: Commit hash không tồn tại $\rightarrow$ HTTP 404. \\\hline
\end{ucspec}

\paragraph{UC-11: Commit qua giao diện web}
\begin{ucspec}{uc11}{Đặc tả Usecase Commit qua giao diện web}
Use Case ID & UC-11 \\\hline
Tên Usecase & Commit qua giao diện web \\\hline
Mô tả & Người dùng chỉnh sửa file đơn lẻ (hoặc nhiều file) trên giao diện web rồi commit. \\\hline
Tác nhân & Người dùng có vai trò $\geq$ WRITE \\\hline
Độ ưu tiên & Cao \\\hline
Tiền điều kiện & Người dùng có quyền ghi vào repo; branch tồn tại. \\\hline
Hậu điều kiện & Commit mới được tạo trên branch chỉ định. \\\hline
Luồng chính &
1. Người dùng mở file trong code browser, nhấn ``Edit'', chỉnh nội dung.\newline
2. Người dùng nhập commit message và nhấn ``Commit changes''.\newline
3. Frontend gọi \texttt{POST /api/v1/repos/:id/commit-file} (hoặc \texttt{/commit-files} cho nhiều file).\newline
4. Core kiểm tra vai trò người dùng trên repo (phải $\geq$ WRITE).\newline
5. Core gọi GitStorage \texttt{CommitFileChange(repoPath, branch, filePath, content, authorName, message)}.\newline
6. GitStorage dùng go-git tạo blob mới, tree mới, commit mới với parent = HEAD branch hiện tại; cập nhật ref.\newline
7. Frontend reload code browser ở branch hiện tại. \\\hline
Luồng ngoại lệ & A1: Không có quyền ghi $\rightarrow$ HTTP 403.\newline A2: File path/Branch không tồn tại $\rightarrow$ HTTP 404. \\\hline
\end{ucspec}

\paragraph{UC-12: So sánh hai branch (compare)}
\begin{ucspec}{uc12}{Đặc tả Usecase So sánh hai branch}
Use Case ID & UC-12 \\\hline
Tên Usecase & So sánh hai branch (compare) \\\hline
Mô tả & Người dùng chọn base và head, hệ thống hiển thị các commit khác biệt và file thay đổi. \\\hline
Tác nhân & Người dùng \\\hline
Độ ưu tiên & Cao \\\hline
Tiền điều kiện & Cả hai ref (base, head) tồn tại trong repo. \\\hline
Hậu điều kiện & Hiển thị danh sách commit, file thay đổi, summary; gợi ý mergeable. \\\hline
Luồng chính &
1. Người dùng vào trang Compare, chọn base (mặc định là default branch) và head.\newline
2. Frontend gọi \texttt{GET /api/v1/repos/:id/compare?base=<b>\&head=<h>}.\newline
3. Core gọi GitStorage \texttt{CompareRefs(repoPath, base, head)}.\newline
4. GitStorage tính: danh sách commit của head không có trong base, danh sách file thay đổi (added/modified/deleted/renamed), summary (commits, files\_changed, additions, deletions, contributor\_count), kiểm tra khả năng merge.\newline
5. Frontend hiển thị tab Commits, tab Files changed, summary; nút ``Create pull request'' nếu \texttt{can\_compare = true}. \\\hline
Luồng ngoại lệ & A1: Base = head $\rightarrow$ thông báo ``There isn't anything to compare''.\newline A2: Ref không tồn tại $\rightarrow$ HTTP 404. \\\hline
\end{ucspec}

\paragraph{UC-13: Tạo Pull Request}
\begin{ucspec}{uc13}{Đặc tả Usecase Tạo Pull Request}
Use Case ID & UC-13 \\\hline
Tên Usecase & Tạo Pull Request \\\hline
Mô tả & Tạo PR từ kết quả compare; PR ở trạng thái OPEN. \\\hline
Tác nhân & Người dùng \\\hline
Độ ưu tiên & Cao \\\hline
Tiền điều kiện & Compare cho thấy có commit khác biệt; source branch khác target branch. \\\hline
Hậu điều kiện & PR được tạo trong \texttt{core.pull\_requests}, có sự kiện \texttt{created} trong \texttt{core.pull\_request\_events}. \\\hline
Luồng chính &
1. Người dùng từ trang Compare nhập title, description, nhấn ``Create pull request''.\newline
2. Frontend gọi \texttt{POST /api/v1/repos/:id/pulls} với body \{title, description, source\_branch, target\_branch\}.\newline
3. Core validate đầu vào (title không rỗng, source $\neq$ target).\newline
4. Core gọi GitStorage \texttt{CompareRefs} để lấy \texttt{source\_commit\_hash} và \texttt{target\_commit\_hash}.\newline
5. Core ghi PR vào \texttt{core.pull\_requests} với trạng thái OPEN.\newline
6. Core ghi sự kiện \texttt{created} vào \texttt{core.pull\_request\_events}.\newline
7. Frontend chuyển tới trang chi tiết PR. \\\hline
Luồng ngoại lệ & A1: Title rỗng $\rightarrow$ HTTP 400.\newline A2: Source = target $\rightarrow$ HTTP 400. \\\hline
\end{ucspec}

\paragraph{UC-15: Merge Pull Request}
\begin{ucspec}{uc15}{Đặc tả Usecase Merge Pull Request}
Use Case ID & UC-15 \\\hline
Tên Usecase & Merge Pull Request \\\hline
Mô tả & Hợp nhất source branch của PR vào target branch. \\\hline
Tác nhân & Người dùng có vai trò $\geq$ MAINTAINER \\\hline
Độ ưu tiên & Cao \\\hline
Tiền điều kiện & PR đang OPEN; source branch không có conflict với target. \\\hline
Hậu điều kiện & Source branch được merge vào target; PR chuyển sang trạng thái MERGED; sự kiện \texttt{merged} được ghi. \\\hline
Luồng chính &
1. Người dùng nhấn nút ``Merge'' ở trang PR.\newline
2. Frontend cho người dùng tuỳ chỉnh commit message (mặc định theo format GitHub).\newline
3. Frontend gọi \texttt{POST /api/v1/repos/:id/pulls/:pull\_id/merge} với body \{commit\_message, commit\_description\}.\newline
4. Core kiểm tra vai trò người dùng ($\geq$ MAINTAINER).\newline
5. Core gọi GitStorage \texttt{CompareRefs} để xác minh khả năng merge.\newline
6. Nếu sạch: gọi \texttt{MergeBranches(repoPath, source, target, mergerName, message)}; cập nhật trạng thái PR thành MERGED; ghi sự kiện.\newline
7. Nếu có xung đột: trả về HTTP 409 với danh sách file conflict; frontend chuyển sang luồng resolve (UC-16). \\\hline
Luồng ngoại lệ & A1: PR đã MERGED/CLOSED $\rightarrow$ HTTP 400.\newline A2: Không có quyền merge $\rightarrow$ HTTP 403.\newline A3: Conflict $\rightarrow$ HTTP 409, chuyển sang UC-16. \\\hline
\end{ucspec}

\paragraph{UC-16: Giải quyết xung đột merge}
\begin{ucspec}{uc16}{Đặc tả Usecase Giải quyết xung đột merge}
Use Case ID & UC-16 \\\hline
Tên Usecase & Giải quyết xung đột (merge conflict) \\\hline
Mô tả & Người dùng chỉnh sửa file xung đột trong giao diện web rồi commit nội dung đã giải quyết. \\\hline
Tác nhân & Người dùng có vai trò $\geq$ MAINTAINER \\\hline
Độ ưu tiên & Cao \\\hline
Tiền điều kiện & Merge ở UC-15 trả về conflict. \\\hline
Hậu điều kiện & Source branch được commit nội dung đã giải quyết; PR có thể merge ở lần thử kế tiếp. \\\hline
Luồng chính &
1. Frontend gọi \texttt{GET /api/v1/repos/:id/pulls/:pull\_id/conflicts} để lấy danh sách file xung đột.\newline
2. Core gọi GitStorage \texttt{GetConflictingFiles}, sau đó \texttt{GetConflictContent} cho từng file để lấy nội dung có chèn marker (\texttt{<<<<<<<}, \texttt{=======}, \texttt{>>>>>>>}).\newline
3. Người dùng chỉnh nội dung file trong trình soạn thảo trên giao diện.\newline
4. Frontend gửi danh sách \texttt{ConflictResolution\{path, resolved\_content\}} qua \texttt{POST .../conflicts/resolve}.\newline
5. Core gọi GitStorage \texttt{ResolveConflictsAndCommit}; GitStorage commit nội dung lên source branch và push.\newline
6. Người dùng quay lại nhấn ``Merge'' ở UC-15, lần này sạch. \\\hline
Luồng ngoại lệ & A1: Resolution thiếu file $\rightarrow$ HTTP 400. \\\hline
\end{ucspec}

\paragraph{UC-17: Revert Pull Request đã merge}
\begin{ucspec}{uc17}{Đặc tả Usecase Revert Pull Request đã merge}
Use Case ID & UC-17 \\\hline
Tên Usecase & Revert Pull Request đã merge \\\hline
Mô tả & Tạo branch revert và commit revert tự động cho PR đã MERGED. \\\hline
Tác nhân & Người dùng có vai trò $\geq$ MAINTAINER \\\hline
Độ ưu tiên & Trung bình \\\hline
Tiền điều kiện & PR ở trạng thái MERGED. \\\hline
Hậu điều kiện & Branch revert được tạo từ target branch của PR gốc, kèm commit revert; người dùng có thể tạo PR revert. \\\hline
Luồng chính &
1. Người dùng nhấn ``Revert'' trên trang PR đã merge.\newline
2. Frontend gọi \texttt{POST /api/v1/repos/:id/pulls/:pull\_id/revert}.\newline
3. Core lấy PR gốc, gọi GitStorage \texttt{CreateRevertBranch(target, revertBranchName, mergeCommitHash, authorName, message)}.\newline
4. GitStorage tạo branch \texttt{revert-<pull\_number>-<source\_branch>} từ target, sinh commit revert (đảo ngược diff của merge commit).\newline
5. Frontend chuyển tới trang Compare để người dùng tạo PR revert. \\\hline
Luồng ngoại lệ & A1: PR chưa MERGED $\rightarrow$ HTTP 400. \\\hline
\end{ucspec}

\paragraph{UC-19: Tạo Issue}
\begin{ucspec}{uc19}{Đặc tả Usecase Tạo Issue}
Use Case ID & UC-19 \\\hline
Tên Usecase & Tạo Issue \\\hline
Mô tả & Người dùng tạo issue mới cho repository. \\\hline
Tác nhân & Người dùng \\\hline
Độ ưu tiên & Cao \\\hline
Tiền điều kiện & Người dùng đã đăng nhập; có quyền truy cập repo. \\\hline
Hậu điều kiện & Issue được tạo ở trạng thái OPEN; sự kiện \texttt{created} được ghi. \\\hline
Luồng chính &
1. Người dùng vào tab Issues, nhấn ``New issue''.\newline
2. Nhập title (bắt buộc), description (Markdown).\newline
3. Tuỳ chọn: gán label, gán assignee.\newline
4. Frontend gọi \texttt{POST /api/v1/repos/:id/issues}.\newline
5. Core ghi vào \texttt{core.issues}, \texttt{core.issue\_events} (\texttt{created}); nếu có label/assignee thì ghi tương ứng vào các bảng phụ.\newline
6. Frontend chuyển tới trang chi tiết issue. \\\hline
Luồng ngoại lệ & A1: Title rỗng $\rightarrow$ HTTP 400. \\\hline
\end{ucspec}

\paragraph{UC-20: Quản lý Issue (đóng/mở/comment)}
\begin{ucspec}{uc20}{Đặc tả Usecase Quản lý Issue}
Use Case ID & UC-20 \\\hline
Tên Usecase & Quản lý Issue \\\hline
Mô tả & Người dùng comment, đóng (kèm lý do)/mở lại issue. \\\hline
Tác nhân & Người dùng \\\hline
Độ ưu tiên & Cao \\\hline
Tiền điều kiện & Issue tồn tại; người dùng có quyền truy cập repo. \\\hline
Hậu điều kiện & Trạng thái/comment được cập nhật; sự kiện được ghi vào \texttt{core.issue\_events}. \\\hline
Luồng chính &
1. \textit{Comment}: người dùng nhập nội dung và nhấn ``Comment''. Frontend gọi \texttt{POST /api/v1/repos/:id/issues/:issue\_id/comments}.\newline
2. \textit{Đóng issue}: chọn lý do COMPLETED/NOT\_PLANNED/DUPLICATE và nhấn ``Close issue''. Frontend gọi \texttt{PATCH .../issues/:issue\_id/status} với \{status: CLOSED, close\_reason: <reason>\}.\newline
3. \textit{Mở lại}: nhấn ``Reopen''. Frontend gọi \texttt{PATCH .../issues/:issue\_id/status} với \{status: OPEN\}.\newline
4. Core validate transition (chỉ cho chuyển nếu khác status hiện tại).\newline
5. Core ghi sự kiện vào \texttt{core.issue\_events}. \\\hline
Luồng ngoại lệ & A1: Đóng issue đã CLOSED $\rightarrow$ HTTP 400. \\\hline
\end{ucspec}

\paragraph{UC-22: Quản lý Collaborator}
\begin{ucspec}{uc22}{Đặc tả Usecase Quản lý Collaborator}
Use Case ID & UC-22 \\\hline
Tên Usecase & Quản lý Collaborator \\\hline
Mô tả & OWNER/MAINTAINER thêm hoặc gỡ thành viên khỏi repo. \\\hline
Tác nhân & OWNER, MAINTAINER \\\hline
Độ ưu tiên & Cao \\\hline
Tiền điều kiện & Người dùng có vai trò $\geq$ MAINTAINER trên repo. \\\hline
Hậu điều kiện & Bản ghi trong \texttt{core.repository\_collaborators} được thêm/gỡ. \\\hline
Luồng chính &
1. \textit{Thêm}: người dùng tìm theo username, chọn vai trò OWNER/MAINTAINER/WRITE. Frontend gọi \texttt{POST /api/v1/repos/:id/collabs}.\newline
2. Core kiểm tra vai trò người gọi và vai trò mục tiêu (chỉ OWNER mới được tạo OWNER).\newline
3. \textit{Gỡ}: nhấn ``Remove'' trên hàng tương ứng. Frontend gọi \texttt{DELETE /api/v1/repos/:id/collabs/:user\_id}.\newline
4. Core ghi/xoá bản ghi tương ứng. \\\hline
Luồng ngoại lệ & A1: Vai trò không hợp lệ $\rightarrow$ HTTP 400.\newline A2: Người gọi không đủ quyền $\rightarrow$ HTTP 403.\newline A3: Tự gỡ OWNER duy nhất $\rightarrow$ HTTP 400. \\\hline
\end{ucspec}

\paragraph{UC-24: Xem Repo Insights}
\begin{ucspec}{uc24}{Đặc tả Usecase Xem Repo Insights}
Use Case ID & UC-24 \\\hline
Tên Usecase & Xem Repo Insights \\\hline
Mô tả & Hiển thị các bảng phân tích cho repository: commit trend 30 ngày, top contributors, branch activity, language breakdown. \\\hline
Tác nhân & Người dùng / Khách (cho repo PUBLIC) \\\hline
Độ ưu tiên & Trung bình \\\hline
Tiền điều kiện & Repo tồn tại; insights snapshot có sẵn (hoặc được tính tự động lần đầu khi tạo repo). \\\hline
Hậu điều kiện & Hiển thị các biểu đồ; người dùng có thể nhấn ``Recompute'' để tính lại. \\\hline
Luồng chính &
1. Người dùng vào tab Insights của repo.\newline
2. Frontend gọi \texttt{GET /api/v1/repos/:id/insights}.\newline
3. Gateway forward sang dịch vụ Insights (do prefix \texttt{/api/v1/insights} hoặc theo cấu hình).\newline
4. Insights service đọc \texttt{insights.repo\_insights} theo \texttt{repo\_id}, trả về snapshot JSON.\newline
5. Frontend vẽ: biểu đồ commit trend (line chart), bảng top contributors, bảng branch activity, biểu đồ tròn language breakdown. \\\hline
Luồng ngoại lệ & A1: Chưa có snapshot $\rightarrow$ trả về null/200, frontend hiện ``Insights are being computed''. \\\hline
\end{ucspec}

\paragraph{UC-25: Recompute Insights}
\begin{ucspec}{uc25}{Đặc tả Usecase Recompute Insights}
Use Case ID & UC-25 \\\hline
Tên Usecase & Recompute Insights \\\hline
Mô tả & Người dùng yêu cầu tính lại insights cho repo. \\\hline
Tác nhân & Người dùng có quyền truy cập repo \\\hline
Độ ưu tiên & Thấp \\\hline
Tiền điều kiện & Repo tồn tại. \\\hline
Hậu điều kiện & \texttt{insights.repo\_insights} được upsert với snapshot mới. \\\hline
Luồng chính &
1. Người dùng nhấn ``Recompute'' trên trang Insights.\newline
2. Frontend gọi \texttt{POST /api/v1/repos/:id/insights/recompute}.\newline
3. Insights service:\newline\hspace*{1em}- Gọi GitStorage \texttt{GetBranches} để lấy danh sách branch.\newline\hspace*{1em}- Gọi \texttt{GetCommits} cho từng branch để lấy lịch sử.\newline\hspace*{1em}- Gọi \texttt{GetLanguageBreakdown} để lấy phân bố ngôn ngữ.\newline
4. Insights tính: commits\_last\_30d, commit\_trend, top\_contributors, branch\_activity, primary\_language.\newline
5. Insights upsert vào \texttt{insights.repo\_insights}.\newline
6. Frontend reload trang Insights. \\\hline
Luồng ngoại lệ & A1: Repo không có commit nào $\rightarrow$ snapshot rỗng. \\\hline
\end{ucspec}

\paragraph{UC-26: Xem Profile Activity}
\begin{ucspec}{uc26}{Đặc tả Usecase Xem Profile Activity}
Use Case ID & UC-26 \\\hline
Tên Usecase & Xem Profile Activity \\\hline
Mô tả & Hiển thị contribution graph 365 ngày kiểu GitHub cùng activity overview. \\\hline
Tác nhân & Người dùng / Khách (cho profile PUBLIC) \\\hline
Độ ưu tiên & Trung bình \\\hline
Tiền điều kiện & User tồn tại. \\\hline
Hậu điều kiện & Hiển thị contribution graph và các thống kê. \\\hline
Luồng chính &
1. Người dùng vào trang \texttt{/<username>}.\newline
2. Frontend gọi \texttt{GET /api/v1/profile/activity?year=<year>}.\newline
3. Gateway forward sang dịch vụ Insights.\newline
4. Insights service:\newline\hspace*{1em}- Lấy danh sách repo của user (đọc \texttt{core.repositories} \& \texttt{core.repository\_collaborators}).\newline\hspace*{1em}- Với mỗi repo, gọi GitStorage \texttt{GetCommits} và lọc theo năm + author = user.\newline\hspace*{1em}- Tính \texttt{contribution\_days} (group by date), \texttt{total\_contributions}, \texttt{top\_repositories}.\newline\hspace*{1em}- Tính \texttt{activity\_chart}: số commit, code reviews, issues, pull requests trong năm (đọc \texttt{core.issues}, \texttt{core.pull\_requests}).\newline
5. Frontend vẽ heatmap 7$\times$53 ô (53 tuần), tô màu theo độ đậm tương ứng số commit. \\\hline
Luồng ngoại lệ & A1: Năm không có dữ liệu $\rightarrow$ heatmap rỗng. \\\hline
\end{ucspec}


\section{Thiết kế kiến trúc hệ thống}

\subsection{Kiến trúc Microservices}
Hệ thống Synergit được thiết kế theo kiến trúc \textit{Microservices} dựa trên triết lý ``Majestic Monolith + Strategic Services'' của GitHub. Hệ thống gồm bốn binary độc lập:

\begin{itemize}
\item \textbf{Gateway (cổng :8080):} Reverse proxy đứng trước, định tuyến request từ frontend tới dịch vụ phù hợp dựa trên prefix đường dẫn. Frontend chỉ cần biết một URL duy nhất.
\item \textbf{Core (cổng :8081 -- ``majestic monolith''):} Chứa phần lớn nghiệp vụ: xác thực (auth), siêu dữ liệu repository, pull request, issue, label, collaborator, star. Sở hữu schema PostgreSQL \texttt{core}.
\item \textbf{GitStorage (cổng :8082):} Dịch vụ duy nhất được phép truy cập hệ thống tệp \texttt{D:\textbackslash SynergitRepo\textbackslash}; phơi các thao tác Git (cây thư mục, nội dung file, log, diff, branch, commit qua web, merge, revert, conflict resolution) qua HTTP. Cũng phơi ba endpoint Git Smart HTTP (\texttt{info/refs}, \texttt{git-upload-pack}, \texttt{git-receive-pack}). Không sở hữu DB schema.
\item \textbf{Insights (cổng :8083):} Dịch vụ phân tích/CPU-bound, lưu kết quả vào schema \texttt{insights.repo\_insights}; được phép \textit{đọc} bảng của \texttt{core} (ngoại lệ duy nhất). Phơi endpoint Repo Insights và Profile Activity.
\end{itemize}

Các dịch vụ giao tiếp với nhau qua HTTP/REST, sử dụng JSON cho payload. Tất cả dịch vụ chia sẻ một biến môi trường \texttt{JWT\_SECRET} để xác thực token cục bộ -- không có dịch vụ ``auth trung tâm''. Sơ đồ kiến trúc tổng thể được trình bày trong Hình~\ref{fig:architecture}.

% ====================== Sơ đồ kiến trúc Microservices Synergit (TikZ) =======
\begin{figure}[H]
\centering
\resizebox{\linewidth}{!}{%
\begin{tikzpicture}[
  font=\sffamily\small,
  tier/.style={draw=sgBlue, very thick, rounded corners, inner sep=10pt},
  box/.style={draw=sgLine, fill=white, rounded corners, align=center,
              minimum width=2.6cm, minimum height=0.85cm, font=\sffamily\footnotesize},
  svc/.style={draw=sgBlue, fill=sgCream, rounded corners, align=center,
              minimum width=3.0cm, minimum height=1.4cm, font=\sffamily\footnotesize, thick},
  gw/.style={draw=sgBlue, fill=sgBlue, text=white, rounded corners, align=center,
              minimum width=4.0cm, minimum height=1.0cm, font=\sffamily\bfseries\footnotesize, thick},
  lbl/.style={font=\sffamily\bfseries\footnotesize, text=sgBlue},
  arr/.style={-{Latex[length=2mm]}, draw=sgBlue, thick},
  arrhttp/.style={-{Latex[length=2mm]}, draw=sgLine, thick, dashed},
]
% ---- Client tier ----
\node[box, fill=sgCream, draw=sgBlue, minimum width=4.6cm] (browser) {Trình duyệt (React SPA)\\\scriptsize VITE\_API\_BASE\_URL = http://localhost:8080};
\begin{scope}[on background layer]
  \node[tier, fill=blue!3, fit=(browser),
        label={[lbl]above:TẦNG CLIENT — React 19 + Vite + TypeScript}] (client) {};
\end{scope}

% ---- Gateway ----
\node[gw, below=1.2cm of browser] (gateway) {Gateway :8080\\\scriptsize Reverse proxy, định tuyến theo prefix};

% ---- Backend services ----
\node[svc, below left=1.7cm and 0.6cm of gateway] (core) {\textbf{Core :8081}\\\scriptsize ``majestic monolith''\\\scriptsize auth · repo metadata\\\scriptsize PR · issue · label · star};
\node[svc, below=1.7cm of gateway] (gitstorage) {\textbf{GitStorage :8082}\\\scriptsize Hệ thống tệp Git\\\scriptsize go-git + Smart HTTP\\\scriptsize tree · log · diff · merge};
\node[svc, below right=1.7cm and 0.6cm of gateway] (insights) {\textbf{Insights :8083}\\\scriptsize Phân tích / CPU-bound\\\scriptsize commit trend · contributors\\\scriptsize language breakdown};

\begin{scope}[on background layer]
  \node[tier, fill=sgBlue!4, fit=(core)(gitstorage)(insights),
        label={[lbl]above:TẦNG BACKEND — Go 1.22 + Gin (4 binary, Clean Architecture, JWT chia sẻ)}] (backend) {};
\end{scope}

% ---- Data tier ----
\node[box, below=1.6cm of core, minimum width=3.0cm] (corepg) {PostgreSQL\\\scriptsize schema: \texttt{core.*}};
\node[box, below=1.6cm of gitstorage, minimum width=3.0cm] (fs) {Hệ thống tệp\\\scriptsize \texttt{D:\textbackslash SynergitRepo\textbackslash}};
\node[box, below=1.6cm of insights, minimum width=3.0cm] (insightspg) {PostgreSQL\\\scriptsize schema: \texttt{insights.*}};

\begin{scope}[on background layer]
  \node[tier, fill=green!4, fit=(corepg)(fs)(insightspg),
        label={[lbl]below:TẦNG DỮ LIỆU — Postgres schema-per-service + Filesystem do GitStorage sở hữu}] (data) {};
\end{scope}

% ---- arrows ----
\draw[arr] (browser) -- (gateway);
\draw[arr] (gateway.south) -- ++(0,-0.4) -| (core.north);
\draw[arr] (gateway.south) -- (gitstorage.north);
\draw[arr] (gateway.south) -- ++(0,-0.4) -| (insights.north);

% Core <-> GitStorage (HTTP, qua port.GitManager)
\draw[arrhttp] (core.east) -- node[above, font=\scriptsize, sloped]{HTTP (port.GitManager)} (gitstorage.west);
% Insights <-> GitStorage
\draw[arrhttp] (insights.west) -- node[above, font=\scriptsize, sloped]{HTTP (read-only Git)} (gitstorage.east);
% Insights -> Core schema (read-only)
\draw[arrhttp, bend right=20] (insights.south) to node[below, font=\scriptsize, sloped]{read-only \texttt{core.*}} (corepg.east);

% Core -> own DB
\draw[arr] (core) -- (corepg);
% GitStorage -> filesystem
\draw[arr] (gitstorage) -- (fs);
% Insights -> own DB
\draw[arr] (insights) -- (insightspg);
\end{tikzpicture}}
\caption{Sơ đồ kiến trúc Microservices của hệ thống Synergit}
\label{fig:architecture}
\end{figure}


\subsection{Vai trò của Clean Architecture trong việc tách dịch vụ}
Mỗi dịch vụ Synergit tuân thủ Clean Architecture với 4 tầng (\textit{domain}, \textit{port}, \textit{usecase}, \textit{adapter}). Điểm đặc biệt: \textbf{interface \texttt{port.GitManager}} (gồm khoảng 25 phương thức như \texttt{InitBareRepo}, \texttt{GetTree}, \texttt{GetCommits}, \texttt{MergeBranches}\ldots) là ``đường nối'' giữa Core và GitStorage.

\begin{itemize}
\item Trong dịch vụ \textbf{GitStorage}, interface \texttt{GitManager} được hiện thực bởi \texttt{LocalGitAdapter} -- adapter thao tác trực tiếp lên hệ thống tệp qua go-git.
\item Trong dịch vụ \textbf{Core} (và Insights), interface \texttt{GitManager} cũng được hiện thực bởi một adapter khác: \texttt{GitStorageHTTPClient} -- adapter chuyển mỗi lời gọi method thành một HTTP request tới GitStorage.
\end{itemize}

\textit{Tầng usecase của Core không biết đến sự khác biệt này} -- nó chỉ biết về interface. Nhờ đó, việc tách GitStorage ra khỏi monolith ban đầu chỉ là việc thay adapter, không cần sửa logic nghiệp vụ.

\subsection{Vòng đời xử lý một request}
Mỗi request HTTP đi qua chuỗi sau:
\begin{enumerate}
\item Browser gọi \texttt{http://gateway:8080/api/v1/...} với header \texttt{Authorization: Bearer <jwt>}.
\item Gateway kiểm tra prefix đường dẫn, chọn upstream, forward request (giữ nguyên header).
\item Dịch vụ đích chạy middleware: ID request, log, CORS, JWT verification (qua secret chia sẻ).
\item Handler bind \& validate request body/query.
\item Handler gọi tầng \textit{usecase} (logic nghiệp vụ thuần).
\item Usecase gọi các \textit{port} (interface) đã được tiêm phụ thuộc:
  \begin{itemize}
  \item \texttt{RepoRepository} $\to$ \texttt{PostgresRepoStore} (đọc/ghi DB).
  \item \texttt{GitManager} $\to$ \texttt{LocalGitAdapter} (trong GitStorage) hoặc \texttt{GitStorageHTTPClient} (trong Core).
  \end{itemize}
\item Kết quả chạy ngược về handler, được serialize thành JSON và trả về client.
\end{enumerate}

\subsection{Luồng tạo và merge Pull Request}
Luồng tạo và merge một pull request là luồng phức tạp nhất, liên quan đến nhiều dịch vụ. Sơ đồ chi tiết được trình bày trong Hình~\ref{fig:prflow}.

% ====================== Luồng tạo & merge Pull Request (TikZ) ==============
\begin{figure}[H]
\centering
\resizebox{0.95\linewidth}{!}{%
\begin{tikzpicture}[
  font=\sffamily\footnotesize, node distance=0.8cm and 2.0cm,
  pstep/.style={draw=sgBlue, fill=sgCream, rounded corners, align=center,
               minimum width=3.6cm, minimum height=0.95cm, font=\sffamily\footnotesize},
  svc/.style={draw=sgLine, fill=white, rounded corners, align=center,
                minimum width=2.8cm, minimum height=0.85cm, font=\sffamily\scriptsize},
  arr/.style={-{Latex[length=2mm]}, draw=sgBlue, thick},
  darr/.style={-{Latex[length=2mm]}, draw=sgLine, thick, dashed},
]
\node[pstep] (s1) {1. Compare hai branch\\\scriptsize Frontend gọi \texttt{GET /repos/:id/compare}};
\node[pstep, below=of s1] (s2) {2. Tạo Pull Request\\\scriptsize \texttt{POST /repos/:id/pulls}\\\scriptsize ghi vào \texttt{pull\_requests} + sự kiện};
\node[pstep, below=of s2] (s3) {3. Merge\\\scriptsize \texttt{POST /repos/:id/pulls/:id/merge}\\\scriptsize Core kiểm tra mergeable};
\node[pstep, below=of s3] (s4a) {4a. Sạch $\rightarrow$ MergeBranches\\\scriptsize Cập nhật \texttt{pull\_requests.status = MERGED}\\\scriptsize Ghi sự kiện \texttt{merged}};
\node[pstep, right=2.0cm of s4a] (s4b) {4b. Xung đột $\rightarrow$ Resolve flow\\\scriptsize \texttt{GET .../conflicts}, người dùng giải quyết\\\scriptsize \texttt{POST .../conflicts/resolve}\\\scriptsize quay lại bước 3};

\node[svc, right=2.0cm of s1] (gw) {Gateway\\\scriptsize :8080};
\node[svc, right=2.0cm of s2] (core) {Core\\\scriptsize :8081};
\node[svc, right=2.0cm of s3] (gs) {GitStorage\\\scriptsize :8082};

\draw[arr] (s1) -- (s2);
\draw[arr] (s2) -- (s3);
\draw[arr] (s3) -- (s4a);
\draw[arr] (s3.east) -- ++(1.2,0) |- (s4b.west);
\draw[darr] (s1.east) -- node[above, font=\scriptsize]{HTTP} (gw.west);
\draw[darr] (s2.east) -- node[above, font=\scriptsize]{HTTP} (core.west);
\draw[darr] (s3.east) -- node[above, font=\scriptsize]{HTTP} (gs.west);
\end{tikzpicture}}
\caption{Luồng tạo và merge Pull Request: từ Compare $\rightarrow$ Tạo PR $\rightarrow$ Merge (sạch hoặc qua resolve conflict)}
\label{fig:prflow}
\end{figure}


\subsection{Sở hữu hệ thống tệp}
Toàn bộ thư mục lưu trữ Git (\texttt{D:\textbackslash SynergitRepo\textbackslash<username>\textbackslash<repo>.git}) chỉ được dịch vụ \textbf{GitStorage} truy cập. Khi Core cần thực hiện một thao tác liên quan đến hệ thống tệp -- ví dụ tạo repository mới, đổi tên thư mục khi đổi username, hoặc bootstrap repo với file README -- Core sẽ gọi qua HTTP tới GitStorage. Cách này đảm bảo:

\begin{itemize}
\item Một nguồn sự thật duy nhất cho dữ liệu Git (không có race condition đọc/ghi từ nhiều dịch vụ).
\item Có thể thay GitStorage bằng triển khai khác (ví dụ chạy trên cluster với Spokes-like replication) mà không cần đụng tới Core.
\item Phù hợp với mô hình của GitHub thực tế (dịch vụ Spokes của họ là chủ duy nhất của object database).
\end{itemize}

\section{Thiết kế cơ sở dữ liệu}

\subsection{Sơ đồ quan hệ (ERD)}
Cơ sở dữ liệu được tổ chức theo \textit{schema-per-service}: dịch vụ Core sở hữu schema \texttt{core.*}, dịch vụ Insights sở hữu schema \texttt{insights.*}. Trong báo cáo này, để dễ theo dõi, ERD được nhóm theo các \textit{bounded context} bên trong schema \texttt{core}: identity, repository catalog, pull request, issue. Quan hệ trong cùng schema (có khoá ngoại) vẽ bằng nét liền; tham chiếu sang schema khác vẽ bằng nét đứt. Chi tiết đầy đủ từng bảng xem mục~\ref{sec:dbdetail}.

\subsubsection{Identity (xác thực và người dùng)}
% ====================== ERD Identity (xác thực và người dùng) ===============
\begin{figure}[H]\centering
\resizebox{0.7\linewidth}{!}{%
\begin{tikzpicture}[node distance=1.2cm and 2.0cm]
\node[erdent, text width=3.6cm] (users) {\textbf{core.users}\\\hrule\vspace{1pt} id (PK)\\ username (unique)\\ email (unique)\\ password\_hash\\ created\_at / updated\_at};
\end{tikzpicture}}
\caption{ERD bounded context Identity -- Bảng \texttt{core.users}}\label{fig:erd-auth}
\end{figure}


\subsubsection{Repository Catalog (siêu dữ liệu repository)}
% ====================== ERD Repository Catalog =============================
\begin{figure}[H]\centering
\resizebox{0.95\linewidth}{!}{%
\begin{tikzpicture}[node distance=1.0cm and 1.6cm]
\node[erdent, text width=3.6cm] (repo) {\textbf{core.repositories}\\\hrule\vspace{1pt} id (PK)\\ name\\ path (unique)\\ description\\ visibility\\ primary\_language};
\node[erdent, right=of repo, text width=3.6cm] (collab) {\textbf{core.repository\_\\collaborators}\\\hrule\vspace{1pt} repository\_id (PK,FK)\\ user\_id (PK,FK)\\ role (OWNER/MAINTAINER/WRITE)\\ created\_at};
\node[erdent, right=of collab, text width=3.0cm] (star) {\textbf{core.repo\_stars}\\\hrule\vspace{1pt} user\_id (PK,FK)\\ repo\_id (PK,FK)\\ created\_at};
\node[erdent, below=of repo, text width=3.4cm] (label) {\textbf{core.labels}\\\hrule\vspace{1pt} id (PK)\\ repo\_id (FK)\\ name\\ color\\ description};
\node[erdext, below=of star] (users) {core.users\\(cùng schema)};

\draw[erdrel] (repo) -- (collab) node[erdcard,midway]{1..n};
\draw[erdrel] (repo) -- (label) node[erdcard,midway]{1..n};
\draw[erdrel] (repo.south east) to[bend left=10] node[erdcard,midway,below]{1..n} (star.south west);
\draw[erdrel] (collab.south) -- node[erdcard,midway]{n..1} (users.north);
\draw[erdrel] (star.south) -- (users.north east);
\end{tikzpicture}}
\caption{ERD bounded context Repository Catalog (\texttt{repositories}, \texttt{collaborators}, \texttt{labels}, \texttt{stars})}\label{fig:erd-repo}
\end{figure}


\subsubsection{Pull Request và sự kiện}
% ====================== ERD Pull Request ===================================
\begin{figure}[H]\centering
\resizebox{0.95\linewidth}{!}{%
\begin{tikzpicture}
\node[erdent, text width=3.8cm] (pr) at (0,0) {\textbf{core.pull\_requests}\\\hrule\vspace{1pt} id (PK)\\ repo\_id (FK)\\ creator\_id (FK)\\ title / description\\ source\_branch / target\_branch\\ source\_commit\_hash / target\_commit\_hash\\ status (OPEN/MERGED/CLOSED)\\ created\_at / updated\_at};
\node[erdent, text width=3.4cm] (ev) at (6.5,2.0) {\textbf{core.pull\_request\_events}\\\hrule\vspace{1pt} id (PK)\\ pull\_request\_id (FK)\\ actor\_id (FK)\\ event\_type\\ created\_at};
\node[erdent, text width=3.4cm] (lbl) at (6.5,-0.2) {\textbf{core.pull\_request\_labels}\\\hrule\vspace{1pt} pull\_request\_id (PK,FK)\\ label\_id (PK,FK)};
\node[erdent, text width=3.4cm] (asn) at (6.5,-2.4) {\textbf{core.pull\_request\_assignees}\\\hrule\vspace{1pt} pull\_request\_id (PK,FK)\\ user\_id (PK,FK)\\ assigned\_at};
\node[erdext, left=1.4cm of pr] (repos) {core.repositories\\(cùng schema)};

\draw[erdrel] (pr.east) to[bend left=10] node[erdcard,near end,above]{1..n} (ev.west);
\draw[erdrel] (pr.east) -- node[erdcard,near end,above]{n..n via labels} (lbl.west);
\draw[erdrel] (pr.east) to[bend right=10] node[erdcard,near end,below]{n..n via users} (asn.west);
\draw[erdrel] (repos.east) -- node[erdcard,midway,above]{1..n} (pr.west);
\end{tikzpicture}}
\caption{ERD bounded context Pull Request (PR + sự kiện + labels + assignees)}\label{fig:erd-pr}
\end{figure}


\subsubsection{Issue, comment và sự kiện}
% ====================== ERD Issue ==========================================
\begin{figure}[H]\centering
\resizebox{0.95\linewidth}{!}{%
\begin{tikzpicture}
\node[erdent, text width=3.8cm] (is) at (0,0) {\textbf{core.issues}\\\hrule\vspace{1pt} id (PK)\\ repo\_id (FK)\\ creator\_id (FK)\\ title / description\\ status (OPEN/CLOSED)\\ close\_reason\\ created\_at / updated\_at};
\node[erdent, text width=3.4cm] (cm) at (6.5,2.2) {\textbf{core.issue\_comments}\\\hrule\vspace{1pt} id (PK)\\ issue\_id (FK)\\ author\_id (FK)\\ body\\ created\_at};
\node[erdent, text width=3.4cm] (ev) at (6.5,0.0) {\textbf{core.issue\_events}\\\hrule\vspace{1pt} id (PK)\\ issue\_id (FK)\\ actor\_id (FK)\\ event\_type\\ created\_at};
\node[erdent, text width=3.4cm] (lb) at (6.5,-2.0) {\textbf{core.issue\_labels}\\\hrule\vspace{1pt} issue\_id (PK,FK)\\ label\_id (PK,FK)};
\node[erdent, text width=3.4cm] (as) at (6.5,-3.7) {\textbf{core.issue\_assignees}\\\hrule\vspace{1pt} issue\_id (PK,FK)\\ user\_id (PK,FK)\\ assigned\_at};
\node[erdext, left=1.4cm of is] (repos) {core.repositories\\(cùng schema)};

\draw[erdrel] (is.east) to[bend left=10] node[erdcard,near end,above]{1..n} (cm.west);
\draw[erdrel] (is.east) -- node[erdcard,near end,above]{1..n} (ev.west);
\draw[erdrel] (is.east) to[bend right=8] node[erdcard,near end,below]{n..n} (lb.west);
\draw[erdrel] (is.east) to[bend right=18] node[erdcard,near end,below]{n..n} (as.west);
\draw[erdrel] (repos.east) -- node[erdcard,midway,above]{1..n} (is.west);
\end{tikzpicture}}
\caption{ERD bounded context Issue (issue + comments + sự kiện + labels + assignees)}\label{fig:erd-issue}
\end{figure}


\subsubsection{Insights (snapshot phân tích)}
% ====================== ERD Insights schema ================================
\begin{figure}[H]\centering
\resizebox{0.85\linewidth}{!}{%
\begin{tikzpicture}[node distance=1.0cm and 2.2cm]
\node[erdent, text width=4.0cm] (ri) {\textbf{insights.repo\_insights}\\\hrule\vspace{1pt} repo\_id (PK)\\ computed\_at\\ commits\_last\_30d\\ commit\_trend (jsonb)\\ top\_contributors (jsonb)\\ branch\_activity (jsonb)\\ primary\_language\\ language\_breakdown (jsonb)\\ last\_error};
\node[erdext, right=of ri] (cr) {core.repositories\\(schema khác,\\tham chiếu mềm)};
\draw[erdsoft] (ri) -- node[erdcard, midway, align=center]{1..1\\soft ref} (cr);
\end{tikzpicture}}
\caption{ERD schema \texttt{insights} (snapshot phân tích cho từng repository, tham chiếu mềm tới schema \texttt{core})}\label{fig:erd-insights}
\end{figure}


\subsection{Danh sách các bảng dữ liệu}
\renewcommand{\arraystretch}{1.2}
\begin{longtable}{|c|>{\ttfamily\small}p{5.0cm}|p{7.0cm}|}
\caption{Danh sách các bảng dữ liệu}\label{tab:tablelist}\\
\hline\rowcolor{sgCream}\textbf{STT} & \normalfont\textbf{Tên bảng} & \textbf{Ý nghĩa}\\\hline\endfirsthead
\hline\rowcolor{sgCream}\textbf{STT} & \normalfont\textbf{Tên bảng} & \textbf{Ý nghĩa}\\\hline\endhead
1  & core.users & Tài khoản người dùng (định danh, email, mật khẩu băm).\\\hline
2  & core.repositories & Siêu dữ liệu repository (tên, đường dẫn, mô tả, visibility, ngôn ngữ chính).\\\hline
3  & core.repository\_collaborators & Quan hệ user--repo với vai trò (OWNER/MAINTAINER/WRITE).\\\hline
4  & core.pull\_requests & Pull request với source/target branch, commit hash, trạng thái.\\\hline
5  & core.pull\_request\_events & Nhật ký sự kiện của PR (created/merged/closed/reopened/commented).\\\hline
6  & core.pull\_request\_labels & Gán label cho PR.\\\hline
7  & core.pull\_request\_assignees & Gán người phụ trách cho PR.\\\hline
8  & core.issues & Issue (title, description, status, close\_reason).\\\hline
9  & core.issue\_assignees & Gán người phụ trách cho issue.\\\hline
10 & core.issue\_events & Nhật ký sự kiện của issue.\\\hline
11 & core.issue\_comments & Bình luận trên issue.\\\hline
12 & core.issue\_labels & Gán label cho issue.\\\hline
13 & core.labels & Label (tên, màu, mô tả) thuộc về một repository.\\\hline
14 & core.repo\_stars & Quan hệ user--repo cho chức năng star.\\\hline
15 & insights.repo\_insights & Snapshot kết quả phân tích cho từng repository.\\\hline
\end{longtable}

\subsection{Thiết kế chi tiết các bảng dữ liệu}\label{sec:dbdetail}
% ====================== Đặc tả chi tiết các bảng dữ liệu =====================

\subsubsection{Bảng core.users}
\begin{dbtable}{db-users}{Chi tiết bảng core.users}
1 & id & uuid & Khoá chính (gen\_random\_uuid).\\\hline
2 & username & varchar(50) & Tên đăng nhập (duy nhất).\\\hline
3 & email & varchar(255) & Email (duy nhất).\\\hline
4 & password\_hash & varchar(255) & Mật khẩu băm bằng bcrypt.\\\hline
5 & created\_at & timestamptz & Thời điểm đăng ký.\\\hline
6 & updated\_at & timestamptz & Thời điểm cập nhật gần nhất (vd. đổi username).\\\hline
\end{dbtable}

\subsubsection{Bảng core.repositories}
\begin{dbtable}{db-repos}{Chi tiết bảng core.repositories}
1 & id & uuid & Khoá chính.\\\hline
2 & name & varchar(255) & Tên repository.\\\hline
3 & path & varchar(1024) & Đường dẫn tương đối tới thư mục repo trên hệ thống tệp, dạng \texttt{<owner>/<name>.git} (duy nhất).\\\hline
4 & created\_at & timestamptz & Thời điểm tạo.\\\hline
5 & description & text & Mô tả repo.\\\hline
6 & visibility & varchar(16) & PUBLIC / PRIVATE.\\\hline
7 & primary\_language & varchar(64) & Ngôn ngữ chính (cập nhật bởi Insights).\\\hline
\end{dbtable}

\subsubsection{Bảng core.repository\_collaborators}
\begin{dbtable}{db-collab}{Chi tiết bảng core.repository\_collaborators}
1 & repository\_id & uuid & Khoá chính (kết hợp), FK $\to$ \texttt{core.repositories(id)} ON DELETE CASCADE.\\\hline
2 & user\_id & uuid & Khoá chính (kết hợp), FK $\to$ \texttt{core.users(id)} ON DELETE CASCADE.\\\hline
3 & role & varchar(20) & Vai trò: OWNER / MAINTAINER / WRITE.\\\hline
4 & created\_at & timestamptz & Thời điểm thêm collaborator.\\\hline
\end{dbtable}

\subsubsection{Bảng core.pull\_requests}
\begin{dbtable}{db-pulls}{Chi tiết bảng core.pull\_requests}
1 & id & uuid & Khoá chính.\\\hline
2 & repo\_id & uuid & FK $\to$ \texttt{core.repositories(id)} ON DELETE CASCADE.\\\hline
3 & creator\_id & uuid & FK $\to$ \texttt{core.users(id)}.\\\hline
4 & title & varchar(255) & Tiêu đề PR.\\\hline
5 & description & text & Mô tả PR (Markdown).\\\hline
6 & source\_branch & varchar(255) & Branch nguồn (head).\\\hline
7 & target\_branch & varchar(255) & Branch đích (base).\\\hline
8 & source\_commit\_hash & varchar(64) & SHA-1 commit head ở thời điểm tạo PR.\\\hline
9 & target\_commit\_hash & varchar(64) & SHA-1 commit base ở thời điểm tạo PR.\\\hline
10 & status & varchar(16) & OPEN / MERGED / CLOSED.\\\hline
11 & created\_at & timestamptz & Thời điểm tạo PR.\\\hline
12 & updated\_at & timestamptz & Thời điểm cập nhật trạng thái cuối.\\\hline
\end{dbtable}

\subsubsection{Bảng core.pull\_request\_events}
\begin{dbtable}{db-prevents}{Chi tiết bảng core.pull\_request\_events}
1 & id & uuid & Khoá chính.\\\hline
2 & pull\_request\_id & uuid & FK $\to$ \texttt{core.pull\_requests(id)} ON DELETE CASCADE.\\\hline
3 & actor\_id & uuid & FK $\to$ \texttt{core.users(id)}.\\\hline
4 & event\_type & varchar(32) & Loại sự kiện: created / merged / closed / reopened / commented \ldots\\\hline
5 & created\_at & timestamptz & Thời điểm sự kiện.\\\hline
\end{dbtable}

\subsubsection{Bảng core.pull\_request\_labels}
\begin{dbtable}{db-prlabel}{Chi tiết bảng core.pull\_request\_labels}
1 & pull\_request\_id & uuid & Khoá chính (kết hợp), FK $\to$ \texttt{core.pull\_requests(id)} ON DELETE CASCADE.\\\hline
2 & label\_id & uuid & Khoá chính (kết hợp), FK $\to$ \texttt{core.labels(id)} ON DELETE CASCADE.\\\hline
\end{dbtable}

\subsubsection{Bảng core.pull\_request\_assignees}
\begin{dbtable}{db-prassign}{Chi tiết bảng core.pull\_request\_assignees}
1 & pull\_request\_id & uuid & Khoá chính (kết hợp), FK $\to$ \texttt{core.pull\_requests(id)} ON DELETE CASCADE.\\\hline
2 & user\_id & uuid & Khoá chính (kết hợp), FK $\to$ \texttt{core.users(id)} ON DELETE CASCADE.\\\hline
3 & assigned\_at & timestamptz & Thời điểm gán.\\\hline
\end{dbtable}

\subsubsection{Bảng core.issues}
\begin{dbtable}{db-issues}{Chi tiết bảng core.issues}
1 & id & uuid & Khoá chính.\\\hline
2 & repo\_id & uuid & FK $\to$ \texttt{core.repositories(id)} ON DELETE CASCADE.\\\hline
3 & creator\_id & uuid & FK $\to$ \texttt{core.users(id)}.\\\hline
4 & title & text & Tiêu đề issue.\\\hline
5 & description & text & Mô tả issue (Markdown).\\\hline
6 & status & varchar(16) & OPEN / CLOSED.\\\hline
7 & close\_reason & varchar(24) & COMPLETED / NOT\_PLANNED / DUPLICATE (NULL nếu OPEN).\\\hline
8 & created\_at & timestamptz & Thời điểm tạo.\\\hline
9 & updated\_at & timestamptz & Thời điểm cập nhật trạng thái cuối.\\\hline
\end{dbtable}

\subsubsection{Bảng core.issue\_assignees}
\begin{dbtable}{db-issueassign}{Chi tiết bảng core.issue\_assignees}
1 & issue\_id & uuid & Khoá chính (kết hợp), FK $\to$ \texttt{core.issues(id)} ON DELETE CASCADE.\\\hline
2 & user\_id & uuid & Khoá chính (kết hợp), FK $\to$ \texttt{core.users(id)} ON DELETE CASCADE.\\\hline
3 & assigned\_at & timestamptz & Thời điểm gán.\\\hline
\end{dbtable}

\subsubsection{Bảng core.issue\_events}
\begin{dbtable}{db-issueevents}{Chi tiết bảng core.issue\_events}
1 & id & uuid & Khoá chính.\\\hline
2 & issue\_id & uuid & FK $\to$ \texttt{core.issues(id)} ON DELETE CASCADE.\\\hline
3 & actor\_id & uuid & FK $\to$ \texttt{core.users(id)}.\\\hline
4 & event\_type & varchar(32) & Loại sự kiện: created / closed / reopened / commented / labeled / assigned \ldots\\\hline
5 & created\_at & timestamptz & Thời điểm sự kiện.\\\hline
\end{dbtable}

\subsubsection{Bảng core.issue\_comments}
\begin{dbtable}{db-issuecmt}{Chi tiết bảng core.issue\_comments}
1 & id & uuid & Khoá chính.\\\hline
2 & issue\_id & uuid & FK $\to$ \texttt{core.issues(id)} ON DELETE CASCADE.\\\hline
3 & author\_id & uuid & FK $\to$ \texttt{core.users(id)}.\\\hline
4 & body & text & Nội dung comment (Markdown).\\\hline
5 & created\_at & timestamptz & Thời điểm comment.\\\hline
\end{dbtable}

\subsubsection{Bảng core.issue\_labels}
\begin{dbtable}{db-issuelabel}{Chi tiết bảng core.issue\_labels}
1 & issue\_id & uuid & Khoá chính (kết hợp), FK $\to$ \texttt{core.issues(id)} ON DELETE CASCADE.\\\hline
2 & label\_id & uuid & Khoá chính (kết hợp), FK $\to$ \texttt{core.labels(id)} ON DELETE CASCADE.\\\hline
\end{dbtable}

\subsubsection{Bảng core.labels}
\begin{dbtable}{db-labels}{Chi tiết bảng core.labels}
1 & id & uuid & Khoá chính.\\\hline
2 & repo\_id & uuid & FK $\to$ \texttt{core.repositories(id)} ON DELETE CASCADE.\\\hline
3 & name & varchar(50) & Tên label (duy nhất trong cùng repo).\\\hline
4 & color & varchar(7) & Mã màu HEX, ví dụ \texttt{\#d73a4a}.\\\hline
5 & description & text & Mô tả label.\\\hline
6 & created\_at & timestamptz & Thời điểm tạo.\\\hline
\end{dbtable}

\subsubsection{Bảng core.repo\_stars}
\begin{dbtable}{db-stars}{Chi tiết bảng core.repo\_stars}
1 & user\_id & uuid & Khoá chính (kết hợp), FK $\to$ \texttt{core.users(id)} ON DELETE CASCADE.\\\hline
2 & repo\_id & uuid & Khoá chính (kết hợp), FK $\to$ \texttt{core.repositories(id)} ON DELETE CASCADE.\\\hline
3 & created\_at & timestamptz & Thời điểm star.\\\hline
\end{dbtable}

\subsubsection{Bảng insights.repo\_insights}
\begin{dbtable}{db-insights}{Chi tiết bảng insights.repo\_insights}
1 & repo\_id & uuid & Khoá chính, tham chiếu mềm tới \texttt{core.repositories(id)}.\\\hline
2 & computed\_at & timestamptz & Thời điểm tính toán snapshot.\\\hline
3 & commits\_last\_30d & int & Số commit trong 30 ngày gần nhất.\\\hline
4 & commit\_trend & jsonb & Mảng \texttt{[\{date, commit\_count\}]} group theo ngày.\\\hline
5 & top\_contributors & jsonb & Mảng \texttt{[\{author\_name, commit\_count\}]} top theo số commit.\\\hline
6 & branch\_activity & jsonb & Mảng \texttt{[\{branch\_name, commit\_count\}]}.\\\hline
7 & primary\_language & varchar(64) & Ngôn ngữ chiếm tỉ lệ cao nhất.\\\hline
8 & language\_breakdown & jsonb & Mảng \texttt{[\{language, bytes, percentage\}]}.\\\hline
9 & last\_error & text & Thông báo lỗi của lần tính cuối (NULL nếu thành công).\\\hline
\end{dbtable}


\section{Thiết kế nghiệp vụ đặc thù}

\subsection{Giao thức Git Smart HTTP}
Synergit hiện thực đầy đủ ba endpoint của Git Smart HTTP ở dịch vụ GitStorage để hỗ trợ \texttt{git clone/push/pull/fetch}:

\begin{itemize}
\item \texttt{GET /:username/:repo.git/info/refs?service=<service>} -- advertise refs.
\item \texttt{POST /:username/:repo.git/git-upload-pack} -- gửi packfile cho client (clone/fetch).
\item \texttt{POST /:username/:repo.git/git-receive-pack} -- nhận packfile từ client (push).
\end{itemize}

Cả ba endpoint được Gateway forward thẳng tới GitStorage. Không yêu cầu xác thực JWT đối với repository PUBLIC; với repository PRIVATE, GitStorage trả về \texttt{401 Unauthorized}.

\subsection{Đảm bảo nhất quán khi đổi username}
Khi người dùng đổi username, có ba thay đổi phải xảy ra cùng lúc:

\begin{enumerate}
\item Cập nhật trường \texttt{username} trong \texttt{core.users}.
\item Đổi tên thư mục \texttt{D:\textbackslash SynergitRepo\textbackslash<old\_username>\textbackslash} $\to$ \texttt{D:\textbackslash SynergitRepo\textbackslash<new\_username>\textbackslash}.
\item Cập nhật cột \texttt{path} trong \texttt{core.repositories} cho mọi repo của user (thay tiền tố \texttt{<old>/} bằng \texttt{<new>/}).
\item Phát hành JWT mới chứa username mới (vì JWT cũ vẫn ghi username cũ).
\end{enumerate}

Dịch vụ Core thực hiện toàn bộ các bước trên trong cùng một handler: cập nhật DB, gọi GitStorage qua HTTP để đổi tên thư mục, cập nhật \texttt{repositories.path}, phát hành JWT mới và trả về client. Frontend nhận JWT mới, lưu vào \texttt{localStorage} và reload trang.

\subsection{Máy trạng thái Pull Request}
Mỗi pull request tuân theo một máy trạng thái (FSM); mọi chuyển trạng thái đều được ghi vào bảng \texttt{core.pull\_request\_events} kèm thời điểm và người thực hiện. Sơ đồ FSM được trình bày trong Hình~\ref{fig:pr-fsm}.

% ====================== Máy trạng thái Pull Request (TikZ) =================
\begin{figure}[H]
\centering
\resizebox{0.92\linewidth}{!}{%
\begin{tikzpicture}[
  font=\sffamily\small, node distance=2.0cm,
  st/.style={draw=sgBlue, fill=sgCream, rounded corners, minimum height=0.95cm,
             minimum width=2.4cm, align=center, font=\sffamily\footnotesize},
  term/.style={st, fill=green!12, draw=sgAccent},
  cancel/.style={st, fill=red!10, draw=red!60},
  arr/.style={-{Latex[length=2mm]}, draw=sgBlue, thick, font=\sffamily\scriptsize},
]
\node[st] (open) {OPEN\\\scriptsize đang mở};
\node[term, right=of open] (merged) {MERGED\\\scriptsize đã merge};
\node[cancel, below=1.4cm of merged] (closed) {CLOSED\\\scriptsize đã đóng};

\draw[arr] (open) -- node[above, midway]{merge} (merged);
\draw[arr] (open) to[bend right=15] node[below, near start]{close} (closed);
\draw[arr] (closed) to[bend right=15] node[above, near end]{reopen} (open);
\draw[arr] (merged) -- node[right, midway]{revert (tạo PR mới)} (closed);
\end{tikzpicture}}
\caption{Máy trạng thái (FSM) vòng đời Pull Request}
\label{fig:pr-fsm}
\end{figure}

\begin{figure}[H]
\centering
\resizebox{0.7\linewidth}{!}{%
\begin{tikzpicture}[
  font=\sffamily\small, node distance=2.0cm,
  st/.style={draw=sgBlue, fill=sgCream, rounded corners, minimum height=0.95cm,
             minimum width=2.4cm, align=center, font=\sffamily\footnotesize},
  cancel/.style={st, fill=red!10, draw=red!60},
  arr/.style={-{Latex[length=2mm]}, draw=sgBlue, thick, font=\sffamily\scriptsize},
]
\node[st] (open) {OPEN\\\scriptsize đang mở};
\node[cancel, right=of open] (closed) {CLOSED\\\scriptsize đã đóng};

\draw[arr] (open) to[bend left=15] node[above, midway]{close (lý do)} (closed);
\draw[arr] (closed) to[bend left=15] node[below, midway]{reopen} (open);
\end{tikzpicture}}
\caption{Máy trạng thái (FSM) vòng đời Issue (lý do đóng: COMPLETED / NOT\_PLANNED / DUPLICATE)}
\label{fig:issue-fsm}
\end{figure}


\subsection{Cơ chế giải quyết xung đột merge}
Khi merge một pull request, dịch vụ Core gọi \texttt{CompareRefs} qua GitStorage để kiểm tra khả năng merge. Nếu \texttt{Mergeable = false}:

\begin{enumerate}
\item Frontend chuyển sang giao diện ``Resolve conflicts'' và gọi \texttt{GET /api/v1/repos/:repo\_id/pulls/:pull\_id/conflicts}.
\item Core gọi GitStorage \texttt{GetConflictingFiles} để lấy danh sách file xung đột; với mỗi file, gọi \texttt{GetConflictContent} để lấy nội dung file đã chèn marker (\texttt{<<<<<<<}, \texttt{=======}, \texttt{>>>>>>>}).
\item Người dùng chỉnh sửa nội dung trên trình soạn thảo trong giao diện web để giải quyết xung đột.
\item Frontend gửi danh sách \texttt{ConflictResolution\{Path, ResolvedContent\}} qua \texttt{POST /api/v1/repos/:repo\_id/pulls/:pull\_id/conflicts/resolve}.
\item Core gọi GitStorage \texttt{ResolveConflictsAndCommit}; GitStorage commit nội dung đã giải quyết lên source branch của PR và push.
\item PR có thể merge ở lần thử tiếp theo.
\end{enumerate}

\subsection{Tính toán Insights theo yêu cầu}
Insights được tính theo phương thức ``recompute on demand'': khi người dùng vào trang Insights của một repo, frontend gọi \texttt{GET /api/v1/repos/:id/insights} để lấy snapshot mới nhất từ \texttt{insights.repo\_insights}. Người dùng có nút ``Recompute'' để chủ động tính lại; khi đó frontend gọi \texttt{POST /api/v1/repos/:id/insights/recompute}, dịch vụ Insights:

\begin{enumerate}
\item Lấy danh sách commit của tất cả branch (qua GitStorage HTTP).
\item Lấy thông tin từng branch (last commit, ahead/behind).
\item Tính \textit{commits\_last\_30d}, \textit{commit\_trend} (group by ngày), \textit{top\_contributors} (group by author).
\item Tính \textit{branch\_activity} (số commit mỗi branch).
\item Gọi GitStorage \texttt{GetLanguageBreakdown} để lấy phân bố ngôn ngữ (theo \% byte).
\item Lưu kết quả vào \texttt{insights.repo\_insights} (upsert theo \texttt{repo\_id}).
\end{enumerate}

Profile Activity tương tự nhưng tổng hợp dữ liệu commit của tất cả repo mà người dùng có quyền, kết hợp với số lượng issue và PR đã tạo từ \texttt{core.*}.
